\papersection{Model}{sec:model}

The model decides which releases to trade (\ref{sec:universe}), turns the price path at entry into a
forecast (\ref{sec:signal}) and turns the forecast into positions (\ref{sec:strategy}). Nothing in it
uses data at or after the entry minute.

\subsection{Event Universe Construction}\label{sec:universe}

Release lines that share at least 80\% of their timestamps are grouped into one \emph{family}, so
each report (CPI, the Employment Situation) is one event. A family is admitted for asset $a$ in year
$y$ if it moves the market more than other releases at the same clock slot. With
$m_T=|p_{T+30}-p_T|$ the absolute 30-minute move after minute $T$,
\begin{equation}\label{eq:impact}
  I_{f,a,y}=\overline{m}_{T\in\mathcal{R}_f}\big/\,\overline{m}_{T\in\mathcal{C}_f},
\end{equation}
where $\mathcal{R}_f$ are the family's earlier releases and $\mathcal{C}_f$ the same clock slot on
days when another family releases there but $f$ does not. A family needs 30 earlier releases and a
Welch $t$-statistic of at least 2 on the log moves; the screen is redone each year from earlier years
only. On \ES\ it admits eight families (payrolls, CPI, ISM Manufacturing and Services, FOMC Minutes,
Retail Sales, Core PCE, Wholesale Inventories). The FOMC decision fails, because at 14:00 its only
control is the Minutes. Each admitted release is traded in five windows; one (asset, release,
window) is a \emph{game}.

\subsection{Signal}\label{sec:signal}

For game $e$ with entry minute $\tau_e$, the path is the 31 log prices ending at entry, standardized,
\begin{equation}\label{eq:path}
  q_{e,i}=(p_{\tau_e-31+i}-\bar p_e)/s_e,\qquad i=1,\ldots,31,
\end{equation}
so DTW compares the shape of the reaction, not its size; the target is the next 30-minute return
$y_e$. Candidates are earlier games on the same asset, clock slot and window, from any family, that
exited before $e$ entered, within three years and at most the 600 most recent (at least 20 are
required). The DTW distance $d_j$ to each candidate uses squared local costs and a Sakoe--Chiba band of
five minutes \citep{Keogh2005}. The $K=15$ nearest neighbours $\mathcal{K}_e$ are weighted by
similarity and age $\alpha_j$ (days), $\omega_j\propto d_j^{-1}e^{-\alpha_j/730}$, and give two
features:
\begin{equation}\label{eq:features}
  \mu_e=\sum_{j\in\mathcal{K}_e}\omega_j y_j,\qquad
  \sigma_e=\Big(\sum_{j\in\mathcal{K}_e}\omega_j (y_j-\mu_e)^2\Big)^{1/2}.
\end{equation}
The mean $\mu_e$ forecasts the direction of the next move; the dispersion $\sigma_e$ measures how much
the analogues disagreed, and forecasts the move's size (\ref{sec:signalquality}). The setting was
chosen on \ES\ in an earlier version of the model and is used unchanged for the other assets.

\subsection{Model}\label{sec:strategy}

\subsubsection{How do we turn this signal into a strategy}

The position in game $e$ is a direction, a gate and a size,
$w_e=\operatorname{sign}(\mu_e)\,G_e\,S_e$, held for 30 minutes from the close of the entry minute.
Games of different families at the same minute are merged into one position. With $n$ ticks of
slippage per side, tick $\delta$ and entry price $P_e$, the net return in bps is
\begin{equation}\label{eq:cost}
  r_e=w_e y_e-|w_e|\cdot 2n\,\delta/P_e\times 10^4 ,\qquad n\in\{0,\tfrac14,\tfrac12\}.
\end{equation}

\subsubsection{Why 5 windows?}\label{sec:five}

Each release is traded in five back-to-back 30-minute windows entering at $T+a_k$,
$a_k\in\{1,31,61,91,121\}$ (Figure~\ref{fig:ladder}). The reaction to a large release lasts more
than 30 minutes; one window gives too few trades to measure a small edge (five give about 225 a year
in \ES); and the windows test where the information is: if the path carries news, the first window
should be best (\ref{sec:hitrate}). No window enters at $T$, whose price predates the release.

\begin{figure}[!tb]
  \centering
  \includegraphics[width=0.48\linewidth]{event_ladder}
  \caption{The anchor ladder: five 30-minute windows after a release at $T$. Each window's path is
  the 31 minutes ending at its entry.}
  \label{fig:ladder}
\end{figure}

\subsubsection{Gating}\label{sec:gating}

The forecast is noisy near zero, so the book trades only its tails,
\begin{equation}\label{eq:gate}
  G_e=\mathbf{1}\{\mu_e\le Q_{0.3}\ \text{or}\ \mu_e\ge Q_{0.7}\},
\end{equation}
with quantiles fixed on the asset's 2013--2020 games: the outer 60\% of forecasts are traded.

\subsubsection{Weighting}

Trades are sized inversely to the dispersion, $S_e=\min(\tilde\sigma/\sigma_e,\,2)$, with
$\tilde\sigma$ its median in the fitting years: consistent analogues trade up to two contracts,
scattered ones less than one. Assets are combined with inverse-volatility weights reset monthly,
\begin{equation}\label{eq:assetweights}
  R_t=\sum_a\lambda_{a,t}R_{a,t},\qquad
  \lambda_{a,t}\propto G^{c}_{a,t}/\hat v_{a,t},
\end{equation}
where $\hat v_{a,t}$ is the 60-session volatility of the futures and $G^c_{a,t}$ a cost gate: an asset
is traded only while a round trip costs less than 10\% of its median 30-minute move on admitted
releases over the past year. \ES\ and \NQ\ always pass (weights about 58\% and 42\%), as does \YM\
(about 36/26/38\% with \ES\ and \NQ); \ZN\ passes only in part of 2023 (\ref{sec:treasury_result}).
